% Einheit 5 von 5 – Rechengesetze und Punkt vor Strich (bank/bruchrechnung/e5.jsonl, zusammenbau v0.1)
%% TODO Zweigzeile Teil 1: was hier gelernt wird (Satz in Schülersprache aus dem Einheitstitel) – Einheit 5
\einheitenkopf[e5]{Einheit 5 von 5 · Rechengesetze und Punkt vor Strich}
\zweigzeile{ab Kl. 6 · keine P10-Aufgabe}
\setcounter{aufgabe}{25}

%% TODO Titel: Ich-kann-Satz fehlt in der Bank (Platzhalter: Kettenname) – Nr. 26
%% TODO Anweisung: ein Satz über der ersten Teilaufgabe fehlt in der Bank – Nr. 26
\begin{aufgabe}{Punkt vor Strich}
%% TODO \rechenplatz{n}: Schrittzahl des Grundfalls fehlt in der Bank (nur bei fester Schreibform, 2.3 b)
\begin{teile}
\teil $7 + 3 \cdot 5$ – kreise die Rechnung ein, die zuerst drankommt.
\teil $5 + 4 \cdot 6$
\teil $2{,}5 + 3 \cdot 1{,}2$
\teil $\frac{1}{6} + \frac{1}{12} \cdot 4$
\teil $(3{,}5 + 1{,}5) \cdot 4$
\teil $0{,}5 \cdot 17 \cdot 2$
\teil Berechne den Wert des Terms $\frac{a + b}{c}$ für $a = 4{,}5$, $b = 1{,}5$ und $c = 4$. \hfill (P10 2021 OS)
\teil Im Schwimmbad kostet der Eintritt für Erwachsene 8,00 €, für Kinder 4,00 €. Montags gibt es $25\,\%$ Rabatt. Familie Berg kommt mit Mutter, Vater, Tochter (9 Jahre) und Oma. Gibt der Term den Rabatt richtig an? \hfill (P10 2015 OS)\\
$\frac{1}{4} \cdot (24\text{ €} + 4\text{ €})$ \janein\\
$(2 \cdot 8\text{ €} + 2 \cdot 4\text{ €}) \cdot \frac{25}{100}$ \janein\\
$25 \cdot (8\text{ €} + 8\text{ €} + 8\text{ €} + 4\text{ €}) : 100$ \janein
\end{teile}
\end{aufgabe}

%% TODO Titel: Ich-kann-Satz fehlt in der Bank (Platzhalter: Kettenname) – Nr. 27
%% TODO Anweisung: ein Satz über der ersten Teilaufgabe fehlt in der Bank – Nr. 27
\begin{aufgabe}{Distributivgesetz (Ausklammern bei Dezimalzahlen)}
\begin{teile}
\teil Rechne geschickt: $3{,}7 \cdot 4 + 6{,}3 \cdot 4$
\end{teile}
\end{aufgabe}

%% TODO Titel: Ich-kann-Satz fehlt in der Bank (Platzhalter: Kettenname) – Nr. 28
%% TODO Anweisung: ein Satz über der ersten Teilaufgabe fehlt in der Bank – Nr. 28
\begin{aufgabe}{Überschlag und Prüfen}
\begin{teile}
\teil Kann $3{,}9 \cdot 5{,}1 = 198{,}9$ stimmen? Prüfe mit Überschlag.
\end{teile}
\end{aufgabe}

%% TODO Titel: Ich-kann-Satz fehlt in der Bank (Platzhalter: Kettenname) – Nr. 29
%% TODO Anweisung: ein Satz über der ersten Teilaufgabe fehlt in der Bank – Nr. 29
\begin{aufgabe}{Punkt vor Strich – Fehler finden, Begründen, Anwendung}
\begin{teile}
\teil Nils rechnet: $6 + 4 \cdot 2{,}5 = 10 \cdot 2{,}5 = 25$. Finde den Fehler.
\teil Warum rechnet man bei $3 + 4 \cdot 5$ zuerst $4 \cdot 5$?
\teil Ben kauft 3 Hefte zu je 1,20 € und 2 Stifte zu je 0,85 €. Er hat 5 € dabei. Reicht das Geld?
\end{teile}
\end{aufgabe}
